% Section 9 --- Discussion (~0.4 pages)
% Topics: based rollups, eager sequencing, A2/A4 for specification, EIP-7702 analogy
%
% REVIEW STATUS: Draft. Based rollup generalization argument (§9.1) and eager
% sequencing treatment (§9.2) are ready for external review.

\section{Discussion}
\label{sec:discussion}

\paragraph{Generalization to Based Rollups.}
Based rollups~\cite{justin-based-rollups} use Ethereum L1 validators as
sequencers rather than a dedicated shared-sequencer layer. The A2/A4 separation
developed in this paper applies directly. A based-rollup sequencer is lazy by
Definition~\ref{def:lazy}: L1 validators include transactions in blocks without
evaluating rollup STFs. The only difference from a dedicated shared sequencer is
that the set of possible sequencers is larger (any L1 validator) and the
commitment is the L1 block itself. The structural condition for PEFO
(\S\ref{sec:pefo-structural}) holds: two based rollups that both use the same
L1 block for inclusion lack execution atomicity across that block boundary.
The based-rollup design space has the same Theorem~\ref{thm:a2-necessary}
escape routes: a based-rollup application seeking A2 must implement capability
(i) or (ii) at the application layer.

\paragraph{Eager Sequencing.}
An \emph{eager} sequencer --- one that executes STFs before committing --- is not
lazy by Definition~\ref{def:lazy} and therefore falls outside the scope of
Theorem~\ref{thm:a2-necessary}. This does not mean eager sequencing is without
issues; it introduces its own complexity (sequencer must maintain full rollup
state, latency increases, MEV extraction shifts to the sequencer). The key
observation is that eagerness is a design choice with non-trivial engineering
cost, and current production shared sequencers (Astria, Espresso, Radius) are
all lazy. An eager sequencer that simulates both STFs before committing provides
capability (i) by construction, but requires access to the full state of both
rollups at commitment time.

\paragraph{A2/A4 Separation as a Specification Tool.}
The A2/A4 separation introduced in \S\ref{sec:taxonomy} has direct value as a
specification tool. When a protocol claims to provide ``atomicity,'' the claim
should be annotated with which level of the A0--A4 hierarchy is targeted. This
is analogous to the ISO~7498 layering convention for networking protocols:
specifying which layer a property operates at prevents conflation in both
design and analysis. We propose that cross-rollup protocol specifications
adopt this convention: state explicitly whether the claimed property is A2
(execution outcome) or A4 (economic externality), and provide the corresponding
closure mechanism.

\paragraph{Relation to EIP-7702.}
EIP-7702 (account code delegation) introduces a new class of cross-context
state interactions: an account's code can be set to a delegate for a single
transaction. This creates a form of within-transaction state coupling that
partially resembles capability (ii): the delegated code can implement an
atomic revert of multiple sub-operations within one L1 transaction. For
cross-rollup applications that include an L1 coordination step, EIP-7702
delegates could provide partial-effect prevention at the L1 level, while the
rollup-level legs remain unprotected. The full interaction between EIP-7702
and cross-rollup A2 is a topic for future work.
